% GENERATED FILE. Edit canonical lessons, then run npm run generate:books.
% canonical-source-hash: fnv1a64:f70a9eb03fa3bc64
% canonical-lessons: ZH-C409-gaofeng, ZH-C409-jiaotong, ZH-W409-tong-through, ZH-C409-tongche, ZH-C409-shoupiaoyuan, ZH-C409-dache

\chapter{Heavy Traffic}
\label{ch:zh-heavy-traffic}
\hlchaptermodality{409}

\begin{chapteropening}
\textbf{By the end of this chapter:} \emph{I can talk about rush hour and the traffic, say a line opens, find the conductor and take a taxi.}

Complete the last lesson of chapter 409: I can talk about rush hour and the traffic, say a line opens, find the conductor and take a taxi.
\end{chapteropening}

\section[gāofēng]{gāofēng — peak (of traffic); rush hour}

\label{lesson:ZH-C409-gaofeng}



\begin{quote}
Before the new one: say the Mandarin for a long-distance coach, then the Mandarin for a starting point.
\end{quote}

\subsection*{You'll want to know: gāofēng}
\emph{gāofēng} — \textquotedblleft{}peak (of traffic); rush hour\textquotedblright{}.

\begin{quote}
\emph{gāofēng de shíhou zài chē lǐ zhǎo bu dào zuòwèi} — At rush hour you can't find a seat on the bus.
\end{quote}

Its characters come later; for now it is a word for the ear.

\begin{cousinweb}
\emph{gāofēng} is built from \zh{高} 'high' and \emph{fēng} 'a peak', both inherited from Middle Chinese, \emph{kaw} and \emph{phjowng} (Guangyun, in Baxter's transcription), which lists the second under another spelling, glossed as a mountain peak.
\end{cousinweb}

\begin{grammarlens}[title={What you've built}]
One more word for getting around.
\end{grammarlens}

\subsection*{Your turn}
\begin{itemize}
\raggedright
  \item \emph{Say it:} \emph{gāofēng}
  \item \emph{Say it:} \emph{gāofēng}, once more
  \item \emph{Say it:} \emph{gāofēng de shíhou zài chē lǐ zhǎo bu dào zuòwèi} — At rush hour you can't find a seat on the bus.
  \item \emph{Recall:} say the Mandarin for a long-distance coach, then the Mandarin for a starting point, then say \emph{gāofēng} again
\end{itemize}

\begin{tcolorbox}[breakable,colback=teal!4,colframe=teal!35!black,title={Before you move on}]
What does gāofēng mean? (\textquotedblleft{}Peak (of traffic); rush hour\textquotedblright{}.) Say it once more.
\end{tcolorbox}
\section[\texorpdfstring{\zh{交通} (jiāotōng)}{交通 (jiāotōng)}]{\zh{交通} (jiāotōng) — traffic; transport}

\label{lesson:ZH-C409-jiaotong}



\begin{quote}
Before the new one: say the Mandarin for a starting point, then the Mandarin for the rush hour.
\end{quote}

\subsection*{You'll want to know: \zh{交通}}
\textbf{\zh{交通}} — \emph{jiāotōng} — \textquotedblleft{}traffic; transport\textquotedblright{}.

\begin{quote}
\emph{zhèr de jiāotōng hěn fāngbiàn} — It's easy to get around here.
\end{quote}

The new character, in pieces: \zh{通} — ten strokes in two pieces: the right piece first, seven strokes, a short bar that turns down to the left with a dot under it, then a frame with two short bars and a vertical through them; the walking piece of \zh{这}, three strokes, sits under it and is written last. The next lesson puts it in your hand.

\begin{cousinweb}
\textbf{\zh{交通}} (\emph{jiāotōng}) is built from \zh{交} 'to cross' and \zh{通} 'to go through', both inherited from Middle Chinese, \emph{kaew} and \emph{thuwng} (Guangyun, in Baxter's transcription), which glosses the second 'to reach through'.
\end{cousinweb}

\begin{grammarlens}[title={What you've built}]
One more word for getting around.
\end{grammarlens}

\subsection*{Your turn}
\begin{itemize}
\raggedright
  \item \emph{Say it:} \emph{jiāotōng}
  \item \emph{Say it:} \emph{jiāotōng}, once more
  \item \emph{Say it:} \emph{zhèr de jiāotōng hěn fāngbiàn} — It's easy to get around here.
  \item \emph{Recall:} say the Mandarin for a starting point, then the Mandarin for the rush hour, then say \emph{jiāotōng} again
\end{itemize}

\begin{tcolorbox}[breakable,colback=teal!4,colframe=teal!35!black,title={Before you move on}]
What does \zh{交通} mean? (\textquotedblleft{}Traffic; transport\textquotedblright{}.) Say it once more.
\end{tcolorbox}
\section[\texorpdfstring{\zh{通} (tōng)}{通 (tōng)}]{\zh{通} — the right piece first, then the walking piece under it}

\label{lesson:ZH-W409-tong-through}



\begin{quote}
The word first, then the character inside it.

\begin{itemize}
\raggedright
  \item \emph{Recall:} say \emph{jiāotōng} — \textbf{R1}, one lesson back, before it decays
\end{itemize}

\begin{quote}
\zh{这} \zh{过}
\end{quote}

\emph{Read it:} each one aloud, with nothing said first — \emph{zhè}, \emph{guò}
\end{quote}

\subsection*{Script — the right piece first, then the walking piece under it}
\hlblockfigure{\detokenize{figures/ZH-W409-tong-through-filmstrip.pdf}}{How \zh{通} is written, stroke by stroke}

\begin{quote}
\zh{通}
\end{quote}

\textbf{Two pieces, and the walking piece comes last.} Count the pieces, not the strokes.

\textbf{The right piece, seven strokes.} A short bar that turns down to the left, and a dot under its end. A vertical down the left. A bar across the top that turns down the right side and hooks to the left at the foot. Two short bars inside. Then a vertical down through the middle.

\textbf{The walking piece, three strokes}, as in \zh{这}: a dot, a stroke that zigzags down, then a long sweep along the foot under the whole character.

7 + 3 = ten.

\subsection*{Your turn}
\begin{itemize}
\raggedright
  \item \emph{Write it:} \zh{通}, all ten, with the model in front of you
  \item \emph{Cover:} write it again from memory, and count the pieces, not the strokes
\end{itemize}

\begin{tcolorbox}[breakable,colback=teal!4,colframe=teal!35!black,title={Before you move on}]
How many strokes in \zh{通}? (\textbf{Ten}.) Which piece comes last? (\textbf{The walking piece, as in \zh{这}}.) What runs down through the middle of the right piece? (\textbf{A vertical}.)

Source: \href{https://raw.githubusercontent.com/chanind/hanzi-writer-data/68d10a4b21150cae5e1ebbd223eed289cf32d90c/data/\%E9\%80\%9A.json}{Hanzi Writer Data \zh{通}, pinned stroke paths and medians}.
\end{tcolorbox}
\section[\texorpdfstring{\zh{通车} (tōngchē)}{通车 (tōngchē)}]{\zh{通车} (tōngchē) — to open to traffic}

\label{lesson:ZH-C409-tongche}



\begin{quote}
Before the new one: say the Mandarin for the rush hour, then the Mandarin for the traffic on the roads.
\end{quote}

\subsection*{You'll want to know: \zh{通车}}
\textbf{\zh{通车}} — \emph{tōngchē} — \textquotedblleft{}to open to traffic\textquotedblright{}.

\begin{quote}
\emph{zhèr de dìtiě míngnián tōngchē} — The metro here opens next year.
\end{quote}

\begin{cousinweb}
\textbf{\zh{通车}} (\emph{tōngchē}) is built from \zh{通} 'to go through' and \zh{车} 'vehicle', both inherited from Middle Chinese, \emph{thuwng} and \emph{tsyhae} (Guangyun, in Baxter's transcription): vehicles can go through.
\end{cousinweb}

\begin{grammarlens}[title={What you've built}]
One more word for getting around.
\end{grammarlens}

\subsection*{Your turn}
\begin{itemize}
\raggedright
  \item \emph{Say it:} \emph{tōngchē}
  \item \emph{Say it:} \emph{tōngchē}, once more
  \item \emph{Say it:} \emph{zhèr de dìtiě míngnián tōngchē} — The metro here opens next year.
  \item \emph{Recall:} say the Mandarin for the rush hour, then the Mandarin for the traffic on the roads, then say \emph{tōngchē} again
\end{itemize}

\begin{tcolorbox}[breakable,colback=teal!4,colframe=teal!35!black,title={Before you move on}]
What does \zh{通车} mean? (\textquotedblleft{}To open to traffic\textquotedblright{}.) Say it once more.
\end{tcolorbox}
\section[shòupiàoyuán]{shòupiàoyuán — ticket seller; bus conductor}

\label{lesson:ZH-C409-shoupiaoyuan}



\begin{quote}
Before the new one: say the Mandarin for the traffic on the roads, then the Mandarin for to open to traffic.
\end{quote}

\subsection*{You'll want to know: shòupiàoyuán}
\emph{shòupiàoyuán} — \textquotedblleft{}ticket seller; bus conductor\textquotedblright{}.

\begin{quote}
\emph{shòupiàoyuán zuò zài ménkǒu} — The conductor sits by the door.
\end{quote}

Its characters come later; for now it is a word for the ear.

\begin{cousinweb}
\emph{shòupiàoyuán} is built from \emph{shòu} 'to sell', \zh{票} 'a ticket' and \emph{yuán} 'a member of staff'. The first and the last are inherited from Middle Chinese, \emph{dzyuwH} and \emph{hjwen} (Guangyun, in Baxter's transcription); no dictionary we use records an older form of \zh{票}, so its history is not settled.
\end{cousinweb}

\begin{grammarlens}[title={What you've built}]
One more word for getting around.
\end{grammarlens}

\subsection*{Your turn}
\begin{itemize}
\raggedright
  \item \emph{Say it:} \emph{shòupiàoyuán}
  \item \emph{Say it:} \emph{shòupiàoyuán}, once more
  \item \emph{Say it:} \emph{shòupiàoyuán zuò zài ménkǒu} — The conductor sits by the door.
  \item \emph{Recall:} say the Mandarin for the traffic on the roads, then the Mandarin for to open to traffic, then say \emph{shòupiàoyuán} again
\end{itemize}

\begin{tcolorbox}[breakable,colback=teal!4,colframe=teal!35!black,title={Before you move on}]
What does shòupiàoyuán mean? (\textquotedblleft{}Ticket seller; bus conductor\textquotedblright{}.) Say it once more.
\end{tcolorbox}
\section[\texorpdfstring{\zh{打车} (dǎchē)}{打车 (dǎchē)}]{\zh{打车} (dǎchē) — to take a taxi}

\label{lesson:ZH-C409-dache}



\begin{quote}
Before the new one: say the Mandarin for to open to traffic, then the Mandarin for a ticket seller.
\end{quote}

\subsection*{You'll want to know: \zh{打车}}
\textbf{\zh{打车}} — \emph{dǎchē} — \textquotedblleft{}to take a taxi\textquotedblright{}.

\begin{quote}
\textbf{\zh{下雨的时候不好打车}} \emph{xià yǔ de shíhou bù hǎo dǎchē} — When it rains it's hard to get a taxi.
\end{quote}

\begin{cousinweb}
\textbf{\zh{打车}} (\emph{dǎchē}) is built from \zh{打} and \zh{车} 'vehicle', both inherited from Middle Chinese, \emph{taengX} and \emph{tsyhae} (Guangyun, in Baxter's transcription), which glosses the first 'to strike'. The sense 'to hire' in this word is not given there, and \emph{taengX} would regularly give a different final, so dǎ is irregular.
\end{cousinweb}

\begin{grammarlens}[title={What you've built}]
One more word for getting around.
\end{grammarlens}

\subsection*{Your turn}
\begin{itemize}
\raggedright
  \item \emph{Say it:} \emph{dǎchē}
  \item \emph{Say it:} \emph{dǎchē}, once more
  \item \emph{Say it:} \emph{xià yǔ de shíhou bù hǎo dǎchē} — When it rains it's hard to get a taxi.
  \item \emph{Recall:} say the Mandarin for to open to traffic, then the Mandarin for a ticket seller, then say \emph{dǎchē} again
\end{itemize}

\begin{tcolorbox}[breakable,colback=teal!4,colframe=teal!35!black,title={Before you move on}]
What does \zh{打车} mean? (\textquotedblleft{}To take a taxi\textquotedblright{}.) Say it once more.
\end{tcolorbox}
